\subsubsection{Challenges}

\paragraph{Optimization}
Maintaining exact equivariance can introduce significant computational overhead.
Models based on irreducible representations are computationally heavy even though recent work has optimized the tensor product decomposition \citep{passaro2023,luo2024a}.
Equivariant architectures may require higher numerical precision, larger memory footprints, longer wall-clock training times relative to unconstrained counterparts and have a dependence to the input size which can limit their scalability to large systems \citep{bogatskiy2020p3icml}.
Optimization difficulties for equivariant neural networks have been observed \citep{pertigkiozoglou2024}.

\paragraph{The ``symmetry bottleneck''}\label{par:sym-bottleneck}
Equivariant neural network methods demand \textit{a priori} knowledge of the data's symmetries before they can be hard coded to the architecture.
While not preventing adoption of these methods, this create a bottleneck because:
\begin{enumerate}
    \item \textbf{Unknown symmetries}: in real datasets symmetries are often unspecified or difficult to define analytically.
    \item \textbf{Manual sophistication}: understanding which symmetry group acts on data or designing a novel architecture requires domain expertise and mathematical sophistication.
    \item \textbf{Inflexibility}: once a symmetry is hard-coded into the architecture, the model cannot adapt if the assumptions are wrong.
    Many real-world datasets only approximately satisfy a given symmetry, for instance due to sensor noise, occlusions, or genuine symmetry breaking in the underlying phenomenon and in that case strictly equivariant models may then be over-constrained
\end{enumerate}
These hurdles are not left unaddressed: (1) by symmetry discovery (\Cref{sec:sym-disc-from-data}); (2) by theoretical addressing multiple groups with the same recipe \citep{ravanbakhsh2017p3icml,finzi2020p3icml,finzi2021p3icml,pearce-crump2023p4icml}, foundational approaches \citep{goel2025} or symmetrization of a backbone \citep{puny2021,kaba2023p4icml}; (3) by approximate soft-equivariance methods (\Cref{sec:soft-equivariance}).

\subsubsection{Equivariant or vanilla architectures ?}
A practical question is whether the design costs, potential difficulties in optimization and hardware regret due to less mature technology of equivariant models is justified over a simpler model, possibly trained with data augmentation.
\citet{nordenfors2024} compare stability of equivariant and vanilla counterpart with data augmentation and \citet{shen2026} provides empirical evidence that, for 3D-rotation equivariant tasks, unconstrained models trained with augmentation reduce their equivariance error rapidly to below 2\% of the total loss within thousands of steps
However equivariant architectures are arguably preferable when data is scarce, when physical consistency of predictions is required, or when generalization to unseen orientations or configurations is critical.
\citet{lawrence2026} propose a decision rule depending on the amount of symmetry bias in the dataset and its correlation with labels and \citet{rygiel2026} find similar conclusion by empirically comparing equivariant and vanilla models on datasets having a varying degree of alignment.
On the theoretical side, \citet{elesedy2021p3icml} obtain a strict generalization benefit in favor of equivariance for linear models.